\chapter{Snoring}

\section*{Definition}
A coarse, harsh breathing sound produced during sleep by vibration of the soft palate, uvula, pharyngeal walls, and tongue base due to partial upper airway obstruction, ranging from benign primary snoring to a marker of obstructive sleep apnea (OSA).

\section*{Pathophysiology}
During sleep, reduced muscle tone in the pharyngeal dilator muscles allows the soft palate, uvula, and tongue to collapse partially into the airway. Airflow through the narrowed oropharynx creates negative pressure causing vibration of floppy soft tissue structures. Factors increasing collapse risk include obesity (fat deposition around airway), supine position, alcohol/sedatives (muscle relaxation), and anatomical narrowing (large tonsils, retrognathia, macroglossia).

\section*{Etiology}
\begin{itemize}
  \item Obesity (most common modifiable risk factor)
  \item Supine sleeping position (gravity worsens airway collapse)
  \item Alcohol, sedatives, or muscle relaxants before sleep
  \item Nasal congestion or obstruction (allergies, deviated septum, polyps)
  \item Enlarged tonsils or adenoids (common in children)
  \item Retrognathia or micrognathia (small or recessed jaw)
  \item Macroglossia (large tongue, seen in Down syndrome, acromegaly)
  \item Hypothyroidism (myxedematous changes, obesity, muscle dysfunction)
  \item Menopause (hormonal changes affecting airway muscle tone)
  \item Aging (decreased pharyngeal muscle tone with age)
\end{itemize}

\section*{Indications for Evaluation}
Seek care if:
\begin{itemize}
  \item Severe or worsening symptoms
  \item Loud snoring with witnessed breathing pauses, choking/gasping during sleep, excessive daytime sleepiness, morning headaches, or hypertension (possible sleep apnea)
  \item Accompanied by other concerning signs
\end{itemize}

\section*{Related Symptoms}
\begin{itemize}
  \item Daytime sleepiness
  \item Fatigue
  \item Morning headache
  \item Sleep apnea
  \item Restless sleep
\end{itemize}
\textbf{Note:} Always consider serious underlying conditions when evaluating persistent snoring.

\section*{Self-Care / Home Management}
\begin{itemize}
  \item Rest and avoid activities that may aggravate the condition.
  \item Maintain adequate hydration and a balanced diet.
  \item Over-the-counter remedies may provide symptomatic relief (consult a pharmacist or doctor).
  \item Monitor symptoms and seek professional advice if they persist or worsen.
  \item Avoid self-medicating with prescription medications without medical guidance.
\end{itemize}

\section*{Diagnostic Approach}
\begin{itemize}
  \item A thorough history and physical examination are the first steps in evaluation.
  \item Laboratory tests (blood work, urinalysis) may be ordered based on clinical suspicion.
  \item Imaging studies (X-ray, ultrasound, CT, MRI) may be indicated when structural pathology is suspected.
  \item Specialist referral is arranged when the diagnosis remains uncertain or requires specific expertise.
\end{itemize}

\section*{Treatment / Management Overview}
\begin{itemize}
  \item Treatment targets the underlying cause identified through diagnostic evaluation.
  \item Supportive care and symptom management are provided while awaiting definitive therapy.
  \item Medications, lifestyle modifications, and physical therapies may be prescribed as appropriate.
  \item Follow-up ensures treatment effectiveness and allows timely adjustments to the care plan.
\end{itemize}

\clearpage